\documentclass[9pt]{article}

\usepackage[utf8]{inputenc}
\usepackage[T1]{fontenc}
\usepackage{mathpazo}
\usepackage{microtype}
\usepackage{needspace}
\usepackage[svgnames]{xcolor}
\usepackage[breaklinks=true,colorlinks=true,linkcolor=deepblue,citecolor=deepblue,urlcolor=deepblue]{hyperref}
\usepackage{geometry}
\usepackage{parskip}
\usepackage{enumitem}

\geometry{a4paper, margin=1.7cm, top=1.3cm, bottom=1.2cm}

% ── Colors ──
\definecolor{deepblue}{RGB}{0, 51, 102}
\definecolor{headinggray}{RGB}{60, 60, 60}
\definecolor{ruleblue}{RGB}{180, 195, 210}

% ── Sections ──
\usepackage{titlesec}
\titleformat{\section}{\large\bfseries\color{deepblue}}{}{0em}{}
\titlespacing{\section}{0pt}{1ex}{0.4ex}

% ── Widow / orphan ──
\clubpenalty=10000
\widowpenalty=10000
\displaywidowpenalty=10000
\raggedbottom

% ── Add color to rules (if booktabs used elsewhere) ──
% (no tables in this proposal that require colored rules)

\begin{document}

% ── Header ──
\noindent{\Huge\bfseries\textcolor{deepblue}{Human--AI Learning Assessment}\par}
\medskip
\noindent{\large\mdseries Testing When Evidence from Student--AI Dialogue Can Be\\[-0.2em]
Attributed to Proficiency Rather Than to Interaction Policy\par}
\medskip
\noindent
\textbf{Proposer:} Fabrizio Terzi, Pyragogy Research\quad
\textbf{Contact:} fabrizio.terzi@pyragogy.org\\
\textbf{Repository:} \url{https://github.com/pyragogy/human-ai-learning-assessment}\quad
\textbf{Date:} October 2026

\medskip
{\color{ruleblue}\rule{\textwidth}{0.7pt}}

\section{Proposed solution and PISA relevance}

\noindent\textbf{The measurement problem.} PISA's transition toward AI-mediated
assessment introduces a validity question that existing frameworks do not answer:
when a student produces a response through dialogue with an AI, can the observed
performance be attributed to the student's proficiency, or does it reflect the
AI's scaffolding, the interaction policy, or the combination? Without a method
to test this attribution, scores from AI-mediated tasks remain ambiguous
regardless of how engaging the interaction is.

\noindent\textbf{Core research programme.} HALA proposes to test \emph{when}
evidence from student--AI dialogue can be attributed to student proficiency
rather than to the interaction policy or the AI's scaffolding. Three
experimental questions structure the investigation:

\begin{enumerate}[leftmargin=1.5em, nosep]
  \item Which student reasoning behaviours are observable and codifiable with
    sufficient reliability across raters and tasks?
  \item How much do these behavioural indicators change when the AI agent,
    its prompting style, and the task structure vary?
  \item Which indicators survive comparison with independent unaided probes
    and tests of cross-linguistic comparability?
\end{enumerate}

\noindent Indicators that fail these tests---that primarily track agent
prompting rather than learner-generated reasoning---will not be advanced as
valid proficiency measures. This is the central falsification criterion of
the programme, not a post-hoc qualification.

\medskip
\noindent\textbf{Illustrative scenario} (hypothetical synthetic stimuli, for
illustration only---no real empirical studies are cited). A student receives
two \emph{synthetic} source cards on screen time and adolescent sleep: a
correlational study reporting a negative association and an experimental
study showing no short-term effect. The AI agent proposes an interpretation
(\textquoteleft Screen time may not affect sleep directly; the correlation could
reflect other factors like evening use\textquoteright). The student can accept,
challenge, request additional evidence, or formulate an alternative explanation.
Two AI roles---explanatory assistant (bounded Socratic prompts) and critical
collaborator (counterexample and verification requests)---enable systematic
comparison across interaction conditions.

\smallskip
\noindent\emph{Worked observable (for the rater manual):} AI turn:\enspace
\textquoteleft Could the correlation reflect a third factor?\textquoteright\enspace Student turn:\enspace
\textquoteleft Maybe kids who use screens at night also sleep less for other reasons\textquoteright.
Rater tags \texttt{COMPETING\_CAUSE}: the student proposes a rival explanation
without asserting it as fact. Failure case: student replies \textquoteleft Yeah,
maybe\textquoteright after a similar AI prompt; the rater records
\texttt{undetermined}---no code fires. The system never marks an utterance
\textquoteleft correct\textquoteright or \textquoteleft incorrect\textquoteright: it records only observables, with
\texttt{learner\_origin\ =\ indeterminate}. Unaided probes interspersed during
the task are a design proposal requiring testing for reactivity, order, and
fatigue---they are not claimed to be causal AI-contribution estimators.\par

\medskip
\noindent\textbf{Measurement development sequence.} Observational codes will
not be presented as a validated instrument at the start of the collaboration.
They will be developed through a predeclared sequence:

\begin{enumerate}[leftmargin=1.5em, nosep]
  \item Define the observation opportunities each task must provide (minimum
    evidence for each code to be scorable);
  \item Annotate the student action and the preceding AI intervention
    separately, preserving turn-level provenance;
  \item Evaluate inter-rater agreement and test whether code distributions
    shift when the AI agent, prompt version, or task changes;
  \item Compare conversational indicators against independent unaided probes
    and examine language comparability;
  \item Retain only indicators that survive these tests; discard those that
    primarily track prompting.
\end{enumerate}

\noindent Candidate codes (draft) include: \textbf{EVIDENCE\_REQUEST}
(student requests grounds for a claim), \textbf{COMPETING\_CAUSE} (proposes
rival explanation), \textbf{DISCRIMINATING\_TEST} (designs comparison to
distinguish hypotheses), \textbf{SOURCE\_APPRAISAL} (evaluates credibility),
\textbf{JUSTIFIED\_REVISION} (revises with evidence link), and
\textbf{CALIBRATED\_UNCERTAINTY} (proportional confidence). Each is
recorded per turn with \texttt{learner\_origin\ =\ indeterminate}---the
system does not identify the cognitive author. A heuristic annotation
pipeline exists (adversarial testing against naive surface scoring) but
produces no proficiency estimate; it is a falsification instrument, not
a scoring system.

\section{Collaboration with OECD}

\noindent\textbf{Proposer delivers:}
\begin{itemize}[leftmargin=1.5em, nosep]
  \item Three task specifications with version-controlled prompt registries
    and a process-coding manual, developed through iterative co-design;
  \item A reproducible synthetic interaction trace pipeline and adversarial
    test suite for detecting prompt-contamination;
  \item A deployment and data-protection design (DPIA-ready framework)
    and a written report on the conditions under which a pilot would
    become feasible.
\end{itemize}

\noindent\textbf{OECD guidance required:}
\begin{itemize}[leftmargin=1.5em, nosep]
  \item Target competency constructs and PISA delivery requirements;
  \item Multilingual adaptation and equivalence standards;
  \item Student safeguarding and data-processing expectations;
  \item Access to authorised pilot partners and psychometric specialists.
\end{itemize}

\noindent\textbf{None of these external inputs or partnerships is secured
today.} The project describes a bounded co-design programme, provisionally
estimated at 4--6 months for the deliverables above and contingent on
collaboration being established and resourcing confirmed; it is not a
committed delivery deadline or SLA. Broader deployment requires partner
agreements and financing; the OECD call does not remunerate the selected
partner.

\section{Intellectual property}

\noindent\textbf{Background:} task concepts, source code, and documentation
developed independently before this collaboration. The repository is private
with no explicit licence; any change to its licensing (for example a
permissive open-source licence) is the owner's decision and remains
conditional on agreement on scope and terms before joint work begins.
\textbf{Foreground:} co-developed specifications and designs. A bounded
research licence to the OECD would be negotiated in writing before work begins.
Third-party rights will be inventoried in the submission.

\section{Privacy, security, and feasibility}

\noindent\textbf{Current state:} the HALA prototype is synthetic only---it
runs on a private VPS, generates no real participant data, and stores no
personal information. Architecture: two-process separation (gateway and
worker) with HMAC-authenticated IPC, audited but uncertified. OS-level
user isolation is not deployed.

\noindent\textbf{Required for any student-facing pilot:} data protection
impact assessment, pseudonymisation, appropriate legal basis under
GDPR/DPA (not automatic parental consent under Art.~8 alone---roles of
controller and processor must be established with the partner institution),
and contractual prohibition on model-provider training from assessment data.

\noindent\textbf{Candidate models.} No model is selected, contracted, or
deployed. Two candidates identified with terms verified against provider
pricing pages on 2026-10-11 (availability may vary by region; re-verified
before any deployment): (1) \textbf{GPT-4o}, released model ID
\texttt{gpt-4o-2024-08-06} (OpenAI API; \$2.50/M input, \$10/M output;
API data are not trained on by default, org-level opt-out detailed at
\url{https://platform.openai.com/docs/pricing}); (2) \textbf{Claude Opus 5.5}
(Anthropic API; \$4/M input, \$20/M output, standard tier; contractual
prohibition on training from submissions at
\url{https://docs.anthropic.com/en/docs/about-claude/pricing}). Self-hosted
alternatives (DeepSeek-V3, Llama~4~Scout) are documented but not assumed
cheaper---at low volume, infrastructure costs may exceed API pricing.
All models require adversarial prompt-leakage testing before assessment use.

\smallskip
\noindent\emph{Cost methodology (indicative, hypothetical).} Configure by token
budget, not by session: turns (\textless 60), input tokens per turn (~800), output
tokens per turn (~300), retries (\textless 5\%)---then multiply by the provider's price
per million tokens. At 50 turns this yields roughly \$0.03--\$0.09 per session
for GPT-4o and \$0.05--\$0.15 for Claude Opus 5.5, before caching or batch
discounts. Final cost depends on prompt structure fixed during development;
no fixed per-session price is claimed.\par

\medskip
{\color{ruleblue}\rule{\textwidth}{0.7pt}}
\medskip

\noindent\textbf{Readiness statement.} This proposal describes a planned
research programme, not an existing operational system. No validated scores,
student data, ethics approval, institutional partnership, or deployed AI
tutor exist. The programme's core question---whether evidence from
AI-mediated dialogue can be validly attributed to student proficiency---is
the object of study, not a capability claim.

\end{document}